
\documentclass[12pt,a4paper]{article}

\usepackage[a4paper,margin=1in]{geometry}
\usepackage{graphicx}
\usepackage{booktabs}
\usepackage{amsmath,amssymb}
\usepackage{float}
\usepackage{hyperref}
\usepackage{xcolor}
\usepackage{listings}
\usepackage{enumitem}
\usepackage{fancyhdr}
\usepackage{longtable}

\hypersetup{colorlinks=true,linkcolor=blue,urlcolor=blue}

\pagestyle{fancy}
\fancyhf{}
\lhead{ICS1512}
\rhead{Machine Learning Algorithms Laboratory}
\cfoot{\thepage}

\lstset{
language=Python,
basicstyle=\ttfamily\small,
numbers=left,
frame=single,
breaklines=true,
showstringspaces=false
}

\begin{document}

\begin{tabular}{|p{4cm}|p{11cm}|}
\hline
Experiment No. & 7 \\ \hline
Experiment Title & Dimensionality Reduction and Model Evaluation (With and Without PCA) \\ \hline

Student Name & Naren Karthik Kandasamy \\ \hline
Register Number & 3122247001036 \\ \hline
Faculty & Dr. Poreddy Ajay Kumar Reddy \\ \hline
Submission Date & 31/8/2026 \\ \hline
\end{tabular}

\section{Objective}
To study the effect of dimensionality reduction using Principal Component Analysis (PCA) on the performance of various machine learning classifiers, comparing their metrics, variance, and bias before and after PCA.

\section{Problem Statement}
We evaluate 10 classification algorithms on the Breast Cancer dataset (30 features) predicting malignancy. We seek to understand whether mapping features to a lower-dimensional principal component space improves efficiency and generalization, and how different algorithms (linear vs ensembles) react to PCA transformations.

\section{Dataset Description}
\begin{longtable}{|p{6cm}|p{8cm}|}
\hline
Dataset Name & Wisconsin Diagnostic Breast Cancer \\ \hline
Dataset Source & UCI Machine Learning Repository \\ \hline
Number of Samples & 569 \\ \hline
Number of Features & 30 \\ \hline
Number of Classes & 2 (Malignant, Benign) \\ \hline
Missing Values & None \\ \hline
Train-Test Split & 80-20 \\ \hline
\end{longtable}

\section{Exploratory Data Analysis}
Include all relevant visualizations.
\begin{itemize}
\item Dataset overview
\item Statistical summary
\item Missing value analysis
\item Class distribution
\item Correlation matrix
\item Feature distribution
\end{itemize}

\subsection*{Figure Template}

\textbf{Inference:}

The original dataset has 30 numerical attributes. Scaling was applied as it is a prerequisite for PCA to prevent high-magnitude features from dominating the principal components.

\section{Data Preprocessing}

Explain:
\begin{itemize}
\item Missing value handling
\item Encoding
\item Feature scaling
\item Train-test split
\item Standardized features using StandardScaler.
\item Applied PCA to retain 95\% of the variance.
\end{itemize}

\section{Machine Learning Algorithm}
Explain:
\begin{itemize}
\item Working principle
\item Mathematical formulation
\item Advantages
\item Limitations
\end{itemize}

\section{Experimental Setup}
\begin{tabular}{ll}
\toprule
Python Version & \\
Scikit-Learn Version & \\
IDE & \\
Operating System & \\
Hardware & \\
\bottomrule
\end{tabular}

\section{Hyperparameter Selection}


\textbf{Table 1: PCA Variance Explained}\\
\begin{tabular}{lllp{4cm}}\toprule Setting & Chosen Components / Variance Target & Explained Variance (\%) & Justification\\ \midrule
With-PCA & 10 Components / 95\% Target & 95.16 & 95\% is a standard threshold to retain max information while shedding noise. \\ \bottomrule \end{tabular}
    
\vspace{1em}\\ 
\textbf{Table 2: SVM Hyperparameter Tuning Results}\\
\begin{tabular}{lll}\toprule Hyperparameters & Performance (No-PCA) & Performance (With-PCA)\\ \midrule
{'C': 0.1, 'kernel': 'linear'} (No-PCA) & 98.25\% & - \\
{'C': 0.1, 'kernel': 'linear'} (With-PCA) & - & 99.12\% \\ \bottomrule \end{tabular}\vspace{1em}\\
        
\textbf{Table 3: Naive Bayes Hyperparameter Tuning Results}\\
\begin{tabular}{lll}\toprule Hyperparameters & Performance (No-PCA) & Performance (With-PCA)\\ \midrule
{'var\_smoothing': 1e-09} (No-PCA) & 96.49\% & - \\
{'var\_smoothing': 1e-09} (With-PCA) & - & 92.11\% \\ \bottomrule \end{tabular}\vspace{1em}\\
        
\textbf{Table 4: KNN Hyperparameter Tuning Results}\\
\begin{tabular}{lll}\toprule Hyperparameters & Performance (No-PCA) & Performance (With-PCA)\\ \midrule
{'n\_neighbors': 5, 'weights': 'uniform'} (No-PCA) & 94.74\% & - \\
{'n\_neighbors': 5, 'weights': 'uniform'} (With-PCA) & - & 95.61\% \\ \bottomrule \end{tabular}\vspace{1em}\\
        
\textbf{Table 5: Logistic Regression Hyperparameter Tuning Results}\\
\begin{tabular}{lll}\toprule Hyperparameters & Performance (No-PCA) & Performance (With-PCA)\\ \midrule
{'C': 1} (No-PCA) & 97.37\% & - \\
{'C': 1} (With-PCA) & - & 98.25\% \\ \bottomrule \end{tabular}\vspace{1em}\\
        
\textbf{Table 6: Decision Tree Hyperparameter Tuning Results}\\
\begin{tabular}{lll}\toprule Hyperparameters & Performance (No-PCA) & Performance (With-PCA)\\ \midrule
{'max\_depth': 3} (No-PCA) & 94.74\% & - \\
{'max\_depth': 5} (With-PCA) & - & 92.98\% \\ \bottomrule \end{tabular}\vspace{1em}\\
        
\textbf{Table 7: Random Forest Hyperparameter Tuning Results}\\
\begin{tabular}{lll}\toprule Hyperparameters & Performance (No-PCA) & Performance (With-PCA)\\ \midrule
{'n\_estimators': 50} (No-PCA) & 96.49\% & - \\
{'n\_estimators': 50} (With-PCA) & - & 94.74\% \\ \bottomrule \end{tabular}\vspace{1em}\\
        
\textbf{Table 8: AdaBoost Hyperparameter Tuning Results}\\
\begin{tabular}{lll}\toprule Hyperparameters & Performance (No-PCA) & Performance (With-PCA)\\ \midrule
{'n\_estimators': 50} (No-PCA) & 96.49\% & - \\
{'n\_estimators': 50} (With-PCA) & - & 95.61\% \\ \bottomrule \end{tabular}\vspace{1em}\\
        
\textbf{Table 9: Gradient Boosting Hyperparameter Tuning Results}\\
\begin{tabular}{lll}\toprule Hyperparameters & Performance (No-PCA) & Performance (With-PCA)\\ \midrule
{'learning\_rate': 0.1, 'n\_estimators': 50} (No-PCA) & 95.61\% & - \\
{'learning\_rate': 0.1, 'n\_estimators': 50} (With-PCA) & - & 96.49\% \\ \bottomrule \end{tabular}\vspace{1em}\\
        
\textbf{Table 10: XGBoost Hyperparameter Tuning Results}\\
\begin{tabular}{lll}\toprule Hyperparameters & Performance (No-PCA) & Performance (With-PCA)\\ \midrule
{'learning\_rate': 0.1, 'n\_estimators': 50} (No-PCA) & 95.61\% & - \\
{'learning\_rate': 0.1, 'n\_estimators': 50} (With-PCA) & - & 96.49\% \\ \bottomrule \end{tabular}\vspace{1em}\\
        
\textbf{Table 11: Stacking Hyperparameter Tuning Results}\\
\begin{tabular}{lll}\toprule Hyperparameters & Performance (No-PCA) & Performance (With-PCA)\\ \midrule
Base: RF, SVM, XGB (No-PCA) & 98.25\% & - \\
Base: RF, SVM, XGB (With-PCA) & - & 97.37\% \\ \bottomrule \end{tabular}\vspace{1em}\\
        
\vspace{1em}\\ 
\textbf{Table 12: 5-Fold Cross-Validation Results (No-PCA vs With-PCA)}\\
\begin{tabular}{llllllll}\toprule Model & Fold 1 & Fold 2 & Fold 3 & Fold 4 & Fold 5 & Avg (No) & Avg (PCA)\\ \midrule
    SVM (No PCA) & 98.2 & 97.4 & 98.2 & 98.2 & 98.2 & 98.07 & - \\\\ SVM (With PCA) & 98.2 & 97.4 & 97.4 & 97.4 & 99.1 & - & 97.89 \\\\ \midrule Naive Bayes (No PCA) & 91.2 & 95.6 & 92.1 & 93.0 & 96.5 & 93.68 & - \\\\ Naive Bayes (With PCA) & 89.5 & 94.7 & 93.0 & 91.2 & 92.0 & - & 92.09 \\\\ \midrule KNN (No PCA) & 98.2 & 97.4 & 97.4 & 93.0 & 94.7 & 96.13 & - \\\\ KNN (With PCA) & 96.5 & 96.5 & 96.5 & 93.9 & 95.6 & - & 95.78 \\\\ \midrule Logistic Regression (No PCA) & 98.2 & 96.5 & 99.1 & 97.4 & 97.3 & 97.71 & - \\\\ Logistic Regression (With PCA) & 99.1 & 96.5 & 98.2 & 98.2 & 98.2 & - & 98.07 \\\\ \midrule Decision Tree (No PCA) & 94.7 & 87.7 & 92.1 & 90.4 & 94.7 & 91.92 & - \\\\ Decision Tree (With PCA) & 96.5 & 94.7 & 94.7 & 93.9 & 92.9 & - & 94.55 \\\\ \midrule Random Forest (No PCA) & 95.6 & 96.5 & 97.4 & 94.7 & 96.5 & 96.13 & - \\\\ Random Forest (With PCA) & 95.6 & 97.4 & 93.0 & 92.1 & 94.7 & - & 94.55 \\\\ \midrule AdaBoost (No PCA) & 98.2 & 97.4 & 99.1 & 93.9 & 96.5 & 97.01 & - \\\\ AdaBoost (With PCA) & 97.4 & 96.5 & 95.6 & 94.7 & 95.6 & - & 95.96 \\\\ \midrule Gradient Boosting (No PCA) & 96.5 & 93.0 & 97.4 & 94.7 & 95.6 & 95.43 & - \\\\ Gradient Boosting (With PCA) & 95.6 & 93.9 & 95.6 & 93.9 & 96.5 & - & 95.08 \\\\ \midrule XGBoost (No PCA) & 98.2 & 94.7 & 96.5 & 94.7 & 95.6 & 95.96 & - \\\\ XGBoost (With PCA) & 98.2 & 95.6 & 95.6 & 95.6 & 96.5 & - & 96.31 \\\\ \midrule Stacking (No PCA) & 98.2 & 96.5 & 99.1 & 98.2 & 98.2 & 98.07 & - \\\\ Stacking (With PCA) & 97.4 & 96.5 & 97.4 & 98.2 & 98.2 & - & 97.54 \\\\ \bottomrule \end{tabular}



\section{Results}

\begin{tabular}{lc}
\toprule
Metric & Value\\
\midrule
See Table 12 above for comprehensive comparison.
\bottomrule
\end{tabular}

\section{Result Visualization}

Include all graphs generated during the experiment.

For \textbf{every graph}, provide:
\begin{itemize}
\item Figure
\item Caption
\item 3--5 line inference
\end{itemize}

Recommended figures:
\begin{enumerate}
\item Correlation Matrix
\item Class Distribution
\item Confusion Matrix
\item ROC Curve
\item Precision-Recall Curve
\item Learning Curve (if applicable)
\item Feature Importance (if applicable)
\end{enumerate}
\begin{figure}[H]
\centering
\includegraphics[width=0.8\linewidth]{Ex7/Scree_Plot.eps}
\caption{Scree Plot: Explained Variance by Components.}
\end{figure}

\section{Discussion}
PCA successfully reduced the feature space from 30 to 10 components while retaining 95\% variance. Tree-based models (Random Forest, XGBoost) generally perform exceptionally well on original feature spaces but can sometimes see a slight drop in accuracy with PCA since orthogonal transformations disrupt individual feature semantics. Linear models (SVM, Logistic Regression) benefited from decorrelated features and reduced dimensions.

\section{Conclusion}
The experiment confirmed that PCA is highly effective for reducing dimensionality without significant loss of information. While PCA improved or maintained the performance of linear and distance-based models (KNN, SVM), ensemble models were slightly less sensitive to the benefits of PCA.

\section{References}
Use IEEE style.

\end{document}
